\begin{thebibliography}{9}

\bibitem{elkan2008}
C.~Elkan and K.~Noto.
\newblock Learning classifiers from only positive and unlabeled data.
\newblock In \emph{Proceedings of the 14th ACM SIGKDD International Conference
  on Knowledge Discovery and Data Mining}, 2008.
\newblock DOI: 10.1145/1401890.1401920.

\bibitem{hosseini2022}
S.~Hosseini, R.~Gottumukkala, S.~Katragadda, R.~T. Bhupatiraju, Z.~Ashkar,
  C.~W. Borst, and K.~Cochran.
\newblock A multimodal sensor dataset for continuous stress detection of nurses
  in a hospital.
\newblock \emph{Scientific Data}, 9:255, 2022.
\newblock DOI: 10.1038/s41597-022-01361-y.

\bibitem{kiryo2017}
R.~Kiryo, G.~Niu, M.~C. du~Plessis, and M.~Sugiyama.
\newblock Positive-unlabeled learning with non-negative risk estimator.
\newblock In \emph{Advances in Neural Information Processing Systems 30}, 2017.
\newblock arXiv:1703.00593.

\bibitem{mathur2024}
A.~Mathur et~al.
\newblock Body sensor-based multimodal nurse stress detection using machine
  learning.
\newblock In \emph{2024 16th International Conference on Communication Systems
  and Networks (COMSNETS)}, 2024.

\bibitem{munoz2015}
M.~L. Munoz, A.~van~Roon, H.~Riese, C.~Thio, E.~Oostenbroek, I.~Westrik,
  E.~J.~C. de~Geus, R.~Gansevoort, J.~Lefrandt, I.~M. Nolte, and H.~Snieder.
\newblock Validity of (ultra-)short recordings for heart rate variability
  measurements.
\newblock \emph{PLOS ONE}, 10(9):e0138921, 2015.
\newblock DOI: 10.1371/journal.pone.0138921.

\bibitem{orini2023}
M.~Orini et~al.
\newblock Long-term association of ultra-short heart rate variability with
  cardiovascular events.
\newblock \emph{Scientific Reports}, 13, 2023.

\end{thebibliography}


\newpage

\appendix

\section{Broader impacts}

The application is monitoring of employees during work, and that is worth
naming plainly. A reliable stress detector deployed in a hospital could direct
support toward staff who need it, which is the motivating case. The same
instrument could as easily support surveillance, scheduling penalties, or
insurance decisions, and the people measured are in a weak position to refuse.
Consent to be monitored by an employer is not freely given in the way consent
to a research study is.

Our result bears on this directly. A detector recovering one episode in ten at
a deployable alarm rate is not fit for any of those uses, including the
benevolent one, and two of ten participants are undetectable at any usable
threshold. We would regard deployment on the strength of results at this level
as harmful, and we note that a published figure of $F_1 = 0.99$ on this same
dataset could be read as licensing exactly that. Reporting the achievable
number honestly is the contribution most relevant to impact here.

The cohort is ten female nurses at one hospital during a pandemic. Any
performance claim generalised beyond that population would be unsupported, and
differential performance across groups cannot be assessed at this sample size.

\section{Compute} All experiments run on a single laptop CPU with no GPU.
Feature extraction over 387 sessions takes 107 s; the complete set
of reported experiments completes in about half an hour beyond feature
extraction. Pretraining over 68k windows takes 192 s for twelve epochs, and the
label-efficiency sweep, which is 120 model fits, takes roughly thirteen
minutes. The audit that precedes
them reads the 3.5 GB archive once per section and is likewise
CPU-bound. No experiment reported here was preceded by a larger unreported
sweep.

\section{Per-participant cohort and negative construction}
\label{app:cohort}

The modelling set behind every reported result: the ten retained
participants at the fixed negative-to-positive ratio of 1.31 adopted in
Section 4.3, for negative-sampling seed 0. A ratio below the 1.31 target
means the eligible negative pool was exhausted before the target was met;
the lowest, at 1.03, sets the residual balance spread of
$1.28\times$ quoted in Table 3.

\begin{table}[h]
\centering
\begin{tabular}{lrrr}
\toprule
Participant & Positive windows & Negative windows & Ratio \\
\midrule
8B & 74 & 76 & 1.03 \\
5C & 68 & 84 & 1.24 \\
94 & 34 & 44 & 1.29 \\
7A & 248 & 324 & 1.31 \\
6B & 182 & 238 & 1.31 \\
E4 & 310 & 406 & 1.31 \\
F5 & 116 & 152 & 1.31 \\
83 & 314 & 412 & 1.31 \\
15 & 79 & 104 & 1.32 \\
BG & 154 & 203 & 1.32 \\
\midrule
Total & 1579 & 2043 & 1.29 \\
\bottomrule
\end{tabular}
\caption{Cohort and negative construction at the adopted ratio of 1.31, sorted by achieved ratio.}
\end{table}

\section{Exclusion cascade}
\label{app:attrition}

Window counts through the exclusion rules, in the order fixed in advance.
Rule numbering follows the analysis plan; rules that remove no windows at
the window level are omitted.

\begin{table}[h]
\centering
\begin{tabular}{cp{5.8cm}rrr}
\toprule
\# & Rule & Entering & Removed & Surviving \\
\midrule
1 & sessions $<$ 5 min, before non-wear removal & 37252 & 32 & 37220 \\
2 & non-wear windows & 37220 & 87 & 37133 \\
3 & dead-EDA sessions & 37133 & 598 & 36535 \\
8 & participant 6D, no eligible negatives & 36535 & 681 & 35854 \\
9 & flat-EDA participants (four) & 35854 & 11520 & 24334 \\
\bottomrule
\end{tabular}
\caption{Exclusion cascade at the window level.}
\end{table}

\section{Per-fold results}
\label{app:perfold}

Leave-one-participant-out folds, averaged over three negative-sampling
seeds, with the across-seed range. Episode recall is at one false alarm
per worn hour. These are the folds behind the aggregate figures in
Section 5; we report them individually because the spread is wide and the
two lowest-count folds are not comparable to the rest.

\begin{table}[h]
\centering
\begin{tabular}{lrrcr}
\toprule
Participant & Episodes & AUC & AUC range & Episode recall \\
\midrule
5C & 4 & 0.388 & 0.365--0.401 & 0.000 \\
BG & 13 & 0.633 & 0.622--0.646 & 0.000 \\
6B & 10 & 0.682 & 0.671--0.701 & 0.200 \\
15 & 10 & 0.811 & 0.794--0.828 & 0.067 \\
E4 & 23 & 0.813 & 0.804--0.825 & 0.232 \\
94 & 5 & 0.835 & 0.816--0.865 & 0.133 \\
83 & 15 & 0.841 & 0.828--0.861 & 0.089 \\
7A & 23 & 0.847 & 0.827--0.867 & 0.029 \\
F5 & 23 & 0.882 & 0.867--0.898 & 0.159 \\
8B & 12 & 0.885 & 0.858--0.909 & 0.028 \\
\midrule
Mean & 138 & 0.762 & & 0.094 \\
\bottomrule
\end{tabular}
\caption{Per-fold AUC and episode recall, seed-averaged. The fold range quoted in the main text, 0.388 to 0.885, is the range of this column.}
\end{table}

\section{Operating points}
\label{app:operating}

The alarm-rate trade-off, averaged over three seeds. The paper reports the
one-alarm-per-worn-hour row; the others are given so the shape of the
curve is visible rather than asserted.

\begin{table}[h]
\centering
\begin{tabular}{rrrrc}
\toprule
False alarms/h & FPR & Window recall & Episode recall & Range \\
\midrule
0.5 & 0.0176 & 0.087 & 0.029 & 0.029--0.029 \\
1.0 & 0.0357 & 0.155 & 0.106 & 0.058--0.130 \\
2.0 & 0.0715 & 0.292 & 0.280 & 0.261--0.319 \\
\bottomrule
\end{tabular}
\caption{Operating points over 138 episodes with usable coverage.}
\end{table}


\newpage
\section*{NeurIPS Paper Checklist}

\begin{enumerate}

\item {\bf Claims}
    \item[] Question: Do the main claims made in the abstract and introduction accurately reflect the paper's contributions and scope?
    \item[] Answer: \answerYes{}
    \item[] Justification: The abstract states the negative result and its boundary explicitly: we eliminated the label-treatment axis and not the representation axis, and Section 7.1 repeats this rather than generalising past it. The intended contribution, an account of the prerequisites for self-supervision on sparsely labelled physiological data, matches what is delivered. Every quantity quoted in the abstract appears with its seed range where one exists.

\item {\bf Limitations}
    \item[] Question: Does the paper discuss the limitations of the work performed by the authors?
    \item[] Answer: \answerYes{}
    \item[] Justification: Section 8 is a dedicated limitations section. It names the narrow cohort, the constructed rather than observed negative class and the consequent inestimability of precision, the unidentifiable class prior, the undertuned positive-unlabelled arm and why its negative result is weaker evidence than the naive comparator, the two-learner search space, the 33 episodes with no sensor coverage, the two normalisation constants chosen rather than derived, the underpowered ten-fold comparisons, and the resolution floor of a twenty-draw permutation test.

\item {\bf Theory assumptions and proofs}
    \item[] Question: For each theoretical result, does the paper provide the full set of assumptions and a complete (and correct) proof?
    \item[] Answer: \answerNA{}
    \item[] Justification: The paper contains no theoretical results. The one formal claim it relies on, that the class prior is not identifiable from positive-unlabelled data, is cited rather than proved.

    \item {\bf Experimental result reproducibility}
    \item[] Question: Does the paper fully disclose all the information needed to reproduce the main experimental results of the paper to the extent that it affects the main claims and/or conclusions of the paper (regardless of whether the code and data are provided or not)?
    \item[] Answer: \answerYes{}
    \item[] Justification: The source dataset is public and cited. The paper specifies the window length and hop, the label rule and its threshold, the exclusion rules and the order in which they are applied, the four eligibility conditions for negatives, the matching procedure and its binning, the negative ratio and how it was chosen, the number of sampling seeds, the feature set by channel, the normalisation and its two constants, the learner, and the fold structure. Where a choice could have gone otherwise, the alternative is reported alongside it rather than omitted.

\item {\bf Open access to data and code}
    \item[] Question: Does the paper provide open access to the data and code, with sufficient instructions to faithfully reproduce the main experimental results, as described in supplemental material?
    \item[] Answer: \answerNo{}
    \item[] Justification: The dataset is public and cited by DOI. The analysis code is not released with this submission, because the repository would identify the authors and the review is double-blind. The protocol is specified in full in Sections 2 through 6 with the intent that it be reimplementable from the paper alone. Code will be released on acceptance or on publication elsewhere.

\item {\bf Experimental setting/details}
    \item[] Question: Does the paper specify all the training and test details (e.g., data splits, hyperparameters, how they were chosen, type of optimizer) necessary to understand the results?
    \item[] Answer: \answerYes{}
    \item[] Justification: Sections 2, 4 and 5 give the splits (leave-one-participant-out, grouped on participant, with the reason stated), the learner and its class weighting, the normalisation window and its floor and clip, the negative ratio and the sensitivity analysis over it, and the class prior grid for the positive-unlabelled arm. Choices fixed before results were seen are identified as such, including the episode-detection fraction and the decision rule in Section 6.

\item {\bf Experiment statistical significance}
    \item[] Question: Does the paper report error bars suitably and correctly defined or other appropriate information about the statistical significance of the experiments?
    \item[] Answer: \answerYes{}
    \item[] Justification: Every quantity computed on the constructed negatives is reported with its range across three negative-sampling seeds, and that range is what the variability captures: the stochastic draw of the comparison group, holding folds, features and learner fixed. The onset analysis reports a sign test and a Wilcoxon signed-rank test rather than a parametric interval, because the underlying distribution is heavy-tailed for reasons the paper explains. The baseline is compared against a twenty-draw within-participant permutation null, reported as mean, standard deviation and maximum. Where the seed range is comparable to the estimate itself, as it is for episode recall, the paper says so at every occurrence rather than only once. Fold-level results are reported individually and never averaged, since two of ten folds behave qualitatively differently from the rest.

\item {\bf Experiments compute resources}
    \item[] Question: For each experiment, does the paper provide sufficient information on the computer resources (type of compute workers, memory, time of execution) needed to reproduce the experiments?
    \item[] Answer: \answerYes{}
    \item[] Justification: A compute section in Appendix B states that all experiments run on a single laptop CPU with no GPU, gives the feature-extraction time over 387 sessions, and states that the complete set of reported experiments finishes in under an hour. It also records that no reported experiment was preceded by a larger unreported sweep.

\item {\bf Code of ethics}
    \item[] Question: Does the research conducted in the paper conform, in every respect, with the NeurIPS Code of Ethics \url{https://neurips.cc/public/EthicsGuidelines}?
    \item[] Answer: \answerYes{}
    \item[] Justification: The work is a secondary analysis of a public, de-identified dataset whose original collection was ethically reviewed by the depositing institution. No new data were collected and no participant was contacted. Anonymity is preserved in this submission.

\item {\bf Broader impacts}
    \item[] Question: Does the paper discuss both potential positive societal impacts and negative societal impacts of the work performed?
    \item[] Answer: \answerYes{}
    \item[] Justification: Appendix A addresses this. The positive case is directing support toward staff who need it. The negative case is that the same instrument supports surveillance, scheduling penalties, or insurance decisions, and that employees are poorly placed to refuse monitoring by an employer. The section argues that a detector performing at the level reported here is unfit for any of these uses including the benevolent one, and notes that a published figure of $F_1 = 0.99$ on this dataset could be read as licensing deployment that our results do not support.

\item {\bf Safeguards}
    \item[] Question: Does the paper describe safeguards that have been put in place for responsible release of data or models that have a high risk for misuse (e.g., pre-trained language models, image generators, or scraped datasets)?
    \item[] Answer: \answerNA{}
    \item[] Justification: No model, dataset, or scraped resource is released. The trained models are not distributed and, on the paper's own argument, are not fit for deployment.

\item {\bf Licenses for existing assets}
    \item[] Question: Are the creators or original owners of assets (e.g., code, data, models), used in the paper, properly credited and are the license and terms of use explicitly mentioned and properly respected?
    \item[] Answer: \answerYes{}
    \item[] Justification: The dataset is credited to its depositors and cited both as a Dryad deposit by DOI and as the accompanying data descriptor. The deposit landing page states that content is licensed for reuse but does not display a specific license designation, so we describe it as publicly available under the repository's terms rather than naming a license we cannot verify. Software used is open source: scikit-learn, NumPy, pandas, and NeuroKit2 for electrodermal response detection.

\item {\bf New assets}
    \item[] Question: Are new assets introduced in the paper well documented and is the documentation provided alongside the assets?
    \item[] Answer: \answerNA{}
    \item[] Justification: No new assets are released with this submission.

\item {\bf Crowdsourcing and research with human subjects}
    \item[] Question: For crowdsourcing experiments and research with human subjects, does the paper include the full text of instructions given to participants and screenshots, if applicable, as well as details about compensation (if any)?
    \item[] Answer: \answerNA{}
    \item[] Justification: This is a secondary analysis of an existing public dataset. We did not recruit, instruct, compensate, or interact with any participant. Instructions given to the original participants and the survey instrument are documented by the depositors in the cited data descriptor.

\item {\bf Institutional review board (IRB) approvals or equivalent for research with human subjects}
    \item[] Question: Does the paper describe potential risks incurred by study participants, whether such risks were disclosed to the subjects, and whether Institutional Review Board (IRB) approvals (or an equivalent approval/review based on the requirements of your country or institution) were obtained?
    \item[] Answer: \answerNA{}
    \item[] Justification: No human subjects research was conducted for this paper. The original collection was reviewed and approved by the depositing institution's review board, as reported in the cited data descriptor; we rely on that record and did not seek separate approval for secondary analysis of de-identified public data.

\item {\bf Declaration of LLM usage}
    \item[] Question: Does the paper describe the usage of LLMs if it is an important, original, or non-standard component of the core methods in this research? Note that if the LLM is used only for writing, editing, or formatting purposes and does \emph{not} impact the core methodology, scientific rigor, or originality of the research, declaration is not required.
    \item[] Answer: \answerNA{}
    \item[] Justification: No large language model is a component of the method. The pipeline is signal processing, feature extraction, and gradient boosting or logistic regression; no LLM runs at training or inference time, and no reported result depends on one. Separately, and as voluntary disclosure rather than because the question requires it: a coding assistant was used while implementing the analysis scripts and drafting the manuscript. Every reported number was produced by those scripts and re-derived from the frozen result tables, and every reference was checked against Crossref, arXiv, or a literature index rather than generated.

\end{enumerate}


\end{document}
